\chapter{Fundamentacao Teorica}
\label{chap_fundamentacao_teorica}

Autômatos celulares sao sistemas dinamicos discretos em espaco, tempo e estado. Em uma malha de \(L\) sitios, a configuracao no instante \(t\) e representada por
\[
\mathbf{x}^{(t)} = \big(x_0^{(t)},x_1^{(t)},\dots,x_{L-1}^{(t)}\big),
\qquad x_i^{(t)}\in\{0,1\}.
\]

Historicamente, esse formalismo foi desenvolvido para investigar auto-organizacao e computacao em sistemas complexos, com bases classicas em von Neumann e desenvolvimentos posteriores em modelagem computacional \cite{vonneumann1966,toffoli1987,chopard1998,ilachinski2001}.

\section{Regra local e dinamica global}

No caso dos automatos celulares elementares, a regra local possui raio \(r=1\) e depende da vizinhanca
\((x_{i-1}^{(t)},x_i^{(t)},x_{i+1}^{(t)})\). A evolucao e dada por
\[
x_i^{(t+1)} = f\big(x_{i-1}^{(t)},x_i^{(t)},x_{i+1}^{(t)}\big),
\]
com condicao de contorno periodica, isto e, os indices sao tomados modulo \(L\).

A funcao local
\[
f:\{0,1\}^3\to\{0,1\}
\]
possui 8 entradas possiveis. Como cada entrada pode ser mapeada para 0 ou 1, existem
\(2^8=256\) regras distintas, numeradas de 0 a 255 \cite{wolfram1983,wolfram2002}.

Equivalentemente, define-se a aplicacao global
\[
F:\{0,1\}^L\to\{0,1\}^L,
\qquad \mathbf{x}^{(t+1)} = F\big(\mathbf{x}^{(t)}\big),
\]
de modo que o automato e um sistema dinamico discreto iterado.

\section{Codificacao numerica da regra de Wolfram}

Se \(a,b,c\in\{0,1\}\), define-se o indice da vizinhanca por
\[
p = 4a + 2b + c \in \{0,\dots,7\}.
\]
Para uma regra \(R\), escreve-se sua forma binaria com 8 bits,
\[
R = \sum_{k=0}^{7} b_k 2^k,
\qquad b_k\in\{0,1\}.
\]
Com a convencao de Wolfram (ordem 111,110,\dots,000), a atualizacao local pode ser escrita como
\[
x_i^{(t+1)} = b_{7-p},
\]
o que implementa diretamente a tabela de transicao da regra.

\section{Metricas matematicas para analise}

Para analisar quantitativamente os padroes gerados, utilizam-se metricas no tempo e no espaco:

\textbf{Densidade de sitios ativos}
\[
\rho(t)=\frac{1}{L}\sum_{i=0}^{L-1} x_i^{(t)}.
\]

\textbf{Atividade temporal}
\[
A(t)=\frac{1}{L}\sum_{i=0}^{L-1}\left|x_i^{(t+1)}-x_i^{(t)}\right|,
\]
que mede a fracao de celulas que mudam de estado entre dois instantes consecutivos.

\textbf{Entropia binaria instantanea}
\[
H(t)=-\rho(t)\log_2\rho(t)-\big(1-\rho(t)\big)\log_2\big(1-\rho(t)\big),
\]
com a convencao de que termos do tipo \(0\log 0\) valem 0.

\textbf{Correlacao espacial de alcance \(r\)}
\[
C_t(r)=\frac{1}{L}\sum_{i=0}^{L-1}x_i^{(t)}x_{i+r}^{(t)}-\rho(t)^2,
\]
util para detectar periodicidade, agrupamento e estrutura de longo alcance.

\textbf{Espalhamento de dano (sensibilidade a perturbacoes)}
\[
d_H(t)=\frac{1}{L}\sum_{i=0}^{L-1}\left|x_i^{(t)}-y_i^{(t)}\right|,
\]
onde \(\mathbf{x}^{(t)}\) e \(\mathbf{y}^{(t)}\) evoluem pela mesma regra a partir de condicoes iniciais quase identicas.

Essas grandezas ajudam a separar regimes dinamicos: regras homogeneas tendem a \(\rho(t)\to 0\) ou 1 e baixa atividade; regras periodicas mantem oscilacoes regulares em \(A(t)\) e \(C_t(r)\); regras caoticas exibem alta entropia e maior crescimento de \(d_H(t)\); e regras complexas combinam ordem local com interacoes nao triviais \cite{wolfram1983,wolfram2002,langton1990,li1990}.

\section{Classes de Wolfram e universalidade}

Os autômatos celulares elementares (Elementary Cellular Automata -- ECA) podem ser modelados como sistemas dinâmicos discretos. Seja o espaço de configurações dado por:

\[
X = \{0,1\}^{\mathbb{Z}}
\]

onde cada célula assume valores binários. A evolução do sistema é definida por uma função global:

\[
x^{(t+1)} = F(x^{(t)})
\]

onde $F$ é induzida por uma regra local:

\[
f : \{0,1\}^3 \rightarrow \{0,1\}
\]

tal que:

\[
(F(x))_i = f(x_{i-1}, x_i, x_{i+1})
\]

A classificação proposta por Wolfram organiza os comportamentos assintóticos desses sistemas em quatro classes fundamentais, com base na evolução a partir de condições iniciais desordenadas.

\subsection{Classe 1 -- Homogênea}

Autômatos da Classe 1 convergem rapidamente para um estado estacionário uniforme. Formalmente, existe um estado fixo $x^*$ tal que:

\[
\exists T \; : \; \forall t \ge T, \quad x^{(t)} = x^*
\]

ou seja, o sistema atinge um ponto fixo do operador $F$:

\[
F(x^*) = x^*
\]

Nesse tipo de comportamento, perturbações nas condições iniciais desaparecem ao longo do tempo, levando o sistema sempre ao mesmo estado final.

\subsection{Classe 2 -- Periódica}

Na Classe 2, o sistema evolui para ciclos periódicos ou estruturas estáveis. Nesse caso, existem inteiros $T$ e $p \geq 1$ tais que:

\[
\forall t \ge T, \quad x^{(t+p)} = x^{(t)}
\]

O sistema entra, portanto, em um ciclo limite. As configurações resultantes consistem de padrões repetitivos ou estruturas locais estáveis que oscilam no tempo.

Perturbações podem persistir, porém permanecem confinadas a regiões limitadas do espaço.

\subsection{Classe 3 -- Caótica}

Autômatos da Classe 3 apresentam comportamento aparentemente aleatório e não periódico. Matematicamente, a órbita do sistema não converge para ciclos de pequeno período e apresenta alta sensibilidade às condições iniciais.

Pequenas perturbações em $x^{(0)}$ resultam em divergências significativas ao longo do tempo:

\[
\| F^t(x) - F^t(x + \delta) \| \not\to 0
\]

Além disso, embora a dinâmica seja determinística, propriedades estatísticas globais (como densidade de células ativas) tendem a valores estáveis.

\subsection{Classe 4 -- Complexa}

A Classe 4 representa um comportamento intermediário entre ordem e caos. Nessa classe, coexistem:

\begin{itemize}
    \item regiões ordenadas (fundos periódicos)
    \item estruturas localizadas persistentes (gliders)
\end{itemize}

Essas estruturas podem se propagar e interagir, formando um sistema altamente dinâmico. O comportamento global não é puramente periódico nem totalmente caótico.

Esse tipo de sistema é particularmente importante, pois pode suportar processamento de informação. Um exemplo clássico é a Regra 110.

\subsection{Universalidade Computacional da Regra 110}

A Regra 110 é um autômato celular elementar que pertence à Classe 4 e apresenta comportamento complexo. Matthew Cook (2004) demonstrou que essa regra é capaz de computação universal, isto é, pode simular uma máquina de Turing.

Isso implica que, em princípio, qualquer algoritmo computacional pode ser emulado por meio da evolução desse autômato. Esse resultado evidencia que regras locais extremamente simples podem gerar comportamentos computacionalmente completos.

\subsection{Discussão}

A classificação de Wolfram permite compreender como diferentes tipos de dinâmica emergem a partir de regras locais simples. Em particular:

\begin{itemize}
    \item Classe 1: comportamento trivial, convergente
    \item Classe 2: comportamento periódico ou estável
    \item Classe 3: comportamento caótico, sensível a condições inicais
    \item Classe 4: comportamento complexo com estruturas emergentes
\end{itemize}

Dessa forma, os autômatos celulares constituem um modelo fundamental para o estudo da emergência de complexidade em sistemas dinâmicos discretos.